% Complete work proof, 19 November 2026; independent review requested.
\begin{result}{F.3.33}{Full profinite cuspidality on every affine hyperbolic curve over a number field}{Proved from the explicit covering maps and F.3.32}
Let $X/k$ be a smooth geometrically connected affine hyperbolic curve over a number field. A nontrivial closed procyclic subgroup of $\pi_1(X_{\overline k})$ has an arithmetic normalizer with open image in $G_k$ exactly when it is contained in a conjugate of an actual cusp inertia subgroup. If
\[
 R_X=\bigcap_\ell\ker\bigl(\pi_1(X_{\overline k})
                 \longrightarrow\pi_1(X_{\overline k})^{(\ell)}\bigr),
\]
then $R_X$ has no nontrivial closed procyclic subgroup with that open-normalizer property.
\end{result}
\begin{proof}

Let $X/k$ be a smooth geometrically connected affine hyperbolic curve
over a number field. Let $\overline X$ be its smooth proper
compactification, of genus $g$, with $n\geq1$ geometric boundary
points. Put $F=\pi_1(X_{\overline k})$ and let
$A=\overline{\langle z\rangle}\subset F$ be a closed procyclic
subgroup whose arithmetic normalizer has open image in $G_k$.
We first give the exact map from a nonperipheral such $z$ to the
single-line witness of F.3.31 and contradict it with F.3.32.
The construction uses the actual covers of the present curve.

\paragraph{A rational cusp and the full-class detector.}
Pass to a finite extension $k'/k$ defining all boundary points.
The intersection of the normalizer with the arithmetic group
over $k'$ still has open image in $G_{k'}$. Choose a boundary
point $c\in\overline X(k')$ and a nonzero vector in its
one-dimensional $k'$-tangent space. The associated actual
tangential section gives
\[
       \Pi=F\rtimes s_c(G_{k'}).
\]
For $z$ outside all conjugate cusp inertia groups,
Nakamura1994, Lemma~2.1.4, printed pp.~96--97\cite{AuditNakamura1994}, applied to
the full class of finite groups, supplies a prime $\ell$
and an open $H\subset F$ containing $A$ such that the
class of $z$ in the proper compactification of the
$H$-cover is nonzero modulo $\ell$. Exactly as in FC.1,
its actual $\ell$-component $u$ has nonzero class in
proper $\Q_\ell$-homology and generates
$A_\ell\simeq\Z_\ell$.

\paragraph{An explicit geometric cover of genus at least three.}
The geometric group $F$ is the profinite completion of the
surface group with presentation
\[
 \langle a_1,b_1,\ldots,a_g,b_g,c_1,\ldots,c_n\mid
        \prod_i[a_i,b_i]\prod_jc_j=1\rangle .
\]
It is free of rank $2g+n-1$. A finite quotient of this
actual geometric group determines a connected finite
etale cover of $X_{\overline k}$. Its compactification
has the indicated inertia indices, so the following
Riemann--Hurwitz calculations are actual covering
calculations.

For $n\geq2$, take the full mod-$5$ abelian quotient
\[
        F\twoheadrightarrow(\Z/5)^{2g+n-1}.
\]
The images of $c_1,\ldots,c_{n-1}$ are independent
boundary coordinates, and the image of $c_n$ is their
negative sum. Each is nonzero, of order~$5$.
Thus the regular cover has degree $d=5^{2g+n-1}$ and
its compactification has genus $g'$ with
\[
       2g'-2=d\left(2g-2+\frac{4n}{5}\right).
                                                               \tag{GA.1}
\]
For $g=0$ hyperbolicity gives $n\geq3$, and the right
side is at least $25\cdot(2/5)=10$. For $g\geq1$
and $n\geq2$, it is still larger. Hence $g'\geq6$.

For $n=1$ and $g\geq2$, take the full mod-$2$
abelian quotient of degree $d=2^{2g}$. The sole
peripheral loop is a product of commutators, hence
has trivial image. The compactification cover is
unramified and
\[
                    g'=1+d(g-1)\geq17.                    \tag{GA.2}
\]

For $(g,n)=(1,1)$, send $a_1,b_1$ to
\[
 x=I+E_{12},\qquad y=I+E_{23}
                  \quad\hbox{in }U_3(\mathbf F_2).
\]
These two elements generate the order-eight upper
unitriangular group. Their commutator is
$I+E_{13}$, of order two. The sole peripheral
loop therefore has inertia order two in this
regular degree-eight cover, giving four boundary
points and
\[
                  2g'-2=8(1-1/2)=4,
                  \qquad g'=3.                            \tag{GA.3}
\]
These cases exhaust all affine hyperbolic $(g,n)$.
Write $H_3\subset F$ for the resulting actual
geometric subgroup. Its finite cover descends to
some finite number field; no chosen descent of
that cover is needed in the next step.

\paragraph{The arithmetic normal core and its proper class.}
Let $S$ be the intersection of all $\Pi$-conjugates
of $H\cap H_3$. The finite generation of $F$ makes
this a finite intersection, so $S$ is open and
$\Pi$-normal. For some $r\geq0$, $u'=u^{\ell^r}$
generates $A_\ell\cap S$. The actual inclusion
$S\hookrightarrow H$ induces proper pushforward
\[
              [u']\longmapsto\ell^r[u]\ne0.               \tag{GA.4}
\]
Thus the proper class on the $S$-cover is nonzero.
Its compactification maps finitely to the
$H_3$-compactification; Riemann--Hurwitz gives
genus at least three.

The subgroup $S\rtimes s_c(G_{k'})$ defines an
actual geometrically connected cover $Y/k'$.
Its intersection with the cusp decomposition
group is
\[
                (S\cap I_c)\rtimes s_c(G_{k'}).
\]
The residue quotient is all of $G_{k'}$, so the
chosen boundary point $y$ of $C=\overline Y$
is rational. Filling all other cusps gives the
actual $k'$-map
\[
                        Y\hookrightarrow C\setminus\{y\}.
                                                               \tag{GA.5}
\]
It preserves the proper homology class (GA.4),
and the image of $A_\ell\cap S$ in the maximal
pro-$\ell$ geometric group of the target is
again $\Z_\ell$, with nonzero degree-one class.
The intersected arithmetic normalizer still
has open image in $G_{k'}$, and every map in
(GA.5) preserves that arithmetic projection.

\paragraph{The actual universal quotient over the same number field.}
Nakamura1996, printed pp.~115--117\cite{AuditNakamura1996}, applies to
curves over an arbitrary number field. Lemma~4.5
makes the actual fibre-image map onto the
universal arithmetic quotient surjective and
records its invariance under finite base
change; the degree-six rank in (4.2)--(4.3)
is nonzero. Choose a finite extension $L/k'$
whose Galois group lies in the open normalizer
image. The same argument as FC.5a/b gives an
actual $D\in U_{g(C),6}$ with
\[
              q_{g(C),6}(D)\ne0,
              \qquad D(a)=[Z_6,a],\qquad a\ne0.           \tag{GA.6}
\]
The exponent is one because the chosen Galois
element fixes the nonzero homology class;
the full-logarithm correction removes the
lower inner terms exactly as in F.3.31.
No rational-descent-to-$\Q$ or CM orbit
argument occurs in (GA.6).

Genus three is covered by the established
single-line statement F.3.29. Every remaining
configuration has $s=g(C)-1\geq3$ and the
actual Levi projection supplies the literal
witness
\[
             M(s)v=0,\qquad LK_{\mathcal D}(s)v\ne0        \tag{GA.7}
\]
of F.3.31. The literal identity of F.3.32 on the seven necessary rows
now gives $LK_{\mathcal D}(s)v=0$, a contradiction. Equivalently,
apply the all-genus single-line vanishing (UC.1) directly to (GA.6).
Thus $z$ belongs to a conjugate of cusp inertia. That inertia group
is closed, so it contains $A=\overline{\langle z\rangle}$.

Conversely, every closed subgroup of a procyclic group is
characteristic. The actual decomposition group of a cusp normalizes
its inertia, and its arithmetic image is the absolute Galois group
of the finite residue extension of that cusp. It therefore has
open image in $G_k$ and normalizes every closed subgroup of the
inertia. This proves the stated equivalence, including before the
finite extension used in the proof.

Finally, the intersection of cusp inertia with $R_X$ is trivial by
an explicit quotient calculation. If $n\geq2$, every peripheral
loop has a nonzero primitive image in the free integral
abelianization: the first $n-1$ are independent coordinates and
the last is their negative sum. Its image in each pro-$\ell$
abelianization consequently has infinite order. If $n=1$,
hyperbolicity gives $g\geq1$ and the peripheral word is
$\prod_{i=1}^g[a_i,b_i]$ up to inversion. Sending $a_1,b_1$ to
$I+E_{12},I+E_{23}$ in $\operatorname{UT}_3(\Z/\ell^m\Z)$
and all the other generators to the identity sends this word to
$I+E_{13}$, of exact order $\ell^m$ for every $m$. Hence in both
cases a cusp exponent $\alpha\in\widehat\Z$ has nontrivial
maximal pro-$\ell$ image whenever its $\ell$-component is
nonzero. Every nonzero exponent has such a component. Therefore
$I\cap R_X=1$, and the same holds for every conjugate inertia
subgroup. The proved cuspidality now excludes a nontrivial subgroup
of $R_X$ with the open-normalizer property.
\end{proof}
